\section{最长完整弦与逐尺度面积支付}\label{sec:chords}
还可以尝试避开预先指定的来源族。对紧致原点星形体 $K$，令 $L_K(q)$ 为 $K$ 中平行于 $q$ 的整条连通线段的最大长度。相应的弦体为
\[
 V(K)=\{A-B:[A,B]\subset K\}.
\]
它一般小于差体 $K-K$：两个端点在 $K$ 内，并不保证连接它们的线段包含于 $K$。它在方向 $q$ 上的径向延伸长度为 $L_K(q)$，由其对称性得到
\[
 |V(K)|=\int_{\mathbb{R}/\pi\mathbb{Z}}L_K(q)^2\,dq.
\]
对满足单位线段可行性的体，$L_K\ge1$。因此，这个内禀量的锐面积比较会很有用。下面的论证给出常数有限的一致比较所需的面积支付，但不确定它的锐常数。

在圆周 $\mathbb R/a\mathbb Z$ 上，记 $\mathcal M f(q)$ 为所有包含 $q$ 的非退化弧 $J$ 上的平均值 $|J|^{-1}\int_J|f|$ 的上确界。这里采用通常的非中心平均极大算子。对 $\alpha>0$，有
\begin{equation}\label{eq:circle-maximal-weak}
 |\{\mathcal M f>\alpha\}|\le\frac3\alpha\|f\|_1.
\end{equation}
证明只用弧覆盖：在一个有限覆盖中，依次选取与已选弧不交的最长弧。每条被弃弧都与某条不短于它的已选弧相交，因此落在后者的三倍扩张内。已选弧不交，且各自积分大于 $\alpha$ 倍弧长，求和即得估计。严格平均不等式允许将弧略微扩大；再对开超水平集取紧集逼近，就得到一般情形。下文取 $a=\pi$，测量的是射影角度长度。

\begin{proposition}[弦长与半径的联合支付]\label{prop:joint-shell}
对原始有限扫掠体 $K$，可测地选取实际最长完整弦，并令 $R(q)$ 为较远端点的模。对二进尺度 $t,M>0$，令
\[
 E_{t,M}=\{q:t\le L_K(q)<2t,\ M\le R(q)<2M\}.
\]
存在有限的绝对常数 $C$，使
\begin{equation}\label{eq:joint-shell}
 Mt\,|E_{t,M}|\le C\,|K\cap\{M/4\le|x|\le M/2\}|.
\end{equation}
因此，对某个有限的绝对常数 $C'$，有 $\int L_K(q)^2\,dq\le C'|K|$。
\end{proposition}

\begin{proof}
有限扫掠体是紧半代数集。关系 $[A,B]\subset K$ 是半代数的，其中保留了对线段参数的全称量词。因此，最大值可以达到，也能用可测规则处理并列情形并选出一条弦。弦长与半径的分箱是 Borel 集。在非空分箱上，三角不等式给出 $t\le4M$。

从较远端点 $A$ 向另一端取长度为 $s=t/32$ 的端部子线段。这条子线段的每一点半径都大于 $M/2$，而两个端点的半径都至多为 $2M$。它与原点张成的三角形，在每个半径 $r\in[M/4,M/2]$ 上都提供两端点射线之间的短角区间。令 $I$ 为这个区间的射影像，$w$ 为其长度。

若支撑高度 $h$ 非零，行列式公式给出
\[
 \sin w=\frac{|h|s}{|A||B'|},\qquad
 w\ge\frac{|h|s}{4M^2},
\]
其中 $B'$ 是端部子线段的另一端点。原弦方向 $q$ 到 $[A]$ 的射影距离至多为一个绝对常数乘以 $|h|/M$。因此，可以选取包含 $q$ 和 $I$ 的区间 $J$，使
\[
 |J|\le w+C_0|h|/M\le C_1(M/t)w.
\]
使用圆上的区间处理射影接缝。

令 $P_r$ 为满足下述条件的射影方向集合：半径为 $r$ 的两个实际点中，至少一个属于 $K$。它的测度至多为半径 $r$ 处实际占据的角度测度。由于 $I\subset P_r$，$1_{P_r}$ 在 $J$ 上的平均值至少为 $c\,t/M$。通常的非中心圆极大不等式将 $E_{t,M}$ 中非零高度部分的测度控制在 $C_2(M/t)|P_r|$ 以内。

在零高度处，端部角区间是一个单点，其射影方向恰好是 $q$。这些方向直接属于 $P_r$。我们单独计入它们的测度：一个点属于 $P_r$，并不保证它对应正的极大平均值。这也说明，不能因为每个角向见证都是零测集，就丢掉连续统多个径向弦。合起来得到
\[
 |E_{t,M}|\le C_3(M/t)|P_r|.
\]
在 $[M/4,M/2]$ 上以实际极坐标因子 $r\,dr$ 积分，就证明了式~\eqref{eq:joint-shell}。

在一个分箱上，$L_K<2t$，因此它对 $\int L_K^2$ 的贡献，至多为一个常数乘以 $t/M$ 再乘以该径向壳层的面积。固定二进尺度 $M$ 后，容许的二进尺度 $t$ 满足 $t\le4M$，而权重 $t/M$ 的几何级数和有一致的界。从同一个壳层支付这整个和。不同二进尺度 $M$ 对应的壳层内部不交，边界圆的面积为零。求和即得最后的断言。
\end{proof}

与第~\ref{sec:radius} 节将使用的远针估计一样，这里由端部子线段产生角区间，再用极大不等式控制来源方向的测度。联合壳层的新增作用在于同时记录弦长和端点半径：先把同一壳层上的几何权重求和，避免为每个弦长类别重复计算整个壳层面积。端部截短、极大估计和求和都留下了定量损失。一个未指定的有限 $C'$ 并不能确定最优比较常数，也不能确定 $A_*$ 的值。

